\section{Related Work}
\label{sec:related}

\paragraph{Morph generation and attack detection.}
Face morphing exploits the possibility that a single enrollment image can resemble multiple contributors, illustrated by the passport attack of Ferrara et al.~\cite{ferrara2014magic}. The SYN-MAD competition evaluates transfer from synthetic training data to real-image attacks~\cite{huber2022synmad}; SMDD established synthetic-only MAD training~\cite{damer2022smdd}. Generation methods differ substantially in how this resemblance is constructed. MIPGAN introduces identity constraints into GAN-based synthesis~\cite{zhang2021mipgan}, whereas MorDIFF uses diffusion autoencoders for representation-level interpolation~\cite{damer2023mordiff}. These mechanisms introduce different reconstruction and processing histories in addition to changes in identity content. Consequently, detecting images produced by a morphing pipeline and detecting evidence attributable to two identities are distinct empirical questions. Our comparisons retain conventional morph-versus-photo evaluation but add controls that hold more of the construction process in common. We do not assume that a shared generator eliminates all low-level differences, nor that an image involving two people necessarily constitutes a successful recognition attack.

\paragraph{Foundation models and synthetic supervision.}
CLIP learns transferable representations through image--text supervision~\cite{radford2021clip}; LoRA adapts pretrained weights using trainable low-rank updates~\cite{hu2022lora}. Their use for morph detection precedes this work. In particular, MADation adapts CLIP with LoRA and a binary classification head~\cite{caldeira2025madation}. Our contribution therefore concerns training examples, controlled evaluation, and the interpretation of detector responses, rather than a new foundation-model architecture. Holding the backbone and optimization recipe fixed allows comparisons of alternative training lists, although removing examples can still change class balance, repetition, and effective diversity. These quantities must be accounted for before attributing an effect specifically to the semantic content of the added controls.

SelfMAD addresses cross-method generalization through self-supervised simulation of morphing artifacts in pixel and frequency space~\cite{ivanovska2025selfmad}. This is closely related to the use of synthetic supervision, but it also highlights a central question for evaluation: which simulated artifacts remain specific to an attack when acquisition and processing change? SelfMAD++ combines such augmentation with a CLIP-based model, local spatial features, and auxiliary segmentation~\cite{ivanovska2026selfmadpp}. We distinguish these two methods; evaluation of the released SelfMAD detector does not establish superiority over SelfMAD++. Our approach is complementary to architecture and augmentation improvements: paired controls can be used to examine their decision mechanisms without presuming that any particular training strategy is inherently artifact-independent.

\paragraph{Shortcut learning and mechanistic evidence.}
Geirhos et al.~\cite{geirhos2020shortcut} describe shortcut learning as the use of predictive rules that succeed under benchmark conditions but fail under more challenging conditions. In morph detection, processing history is a plausible source of such prediction because the two classes can differ systematically in resampling, storage, reconstruction, or image boundaries. A strong benchmark result alone cannot distinguish these explanations from sensitivity to identity mixing. Conversely, a processing perturbation that changes a detector's output does not prove that processing explains every correct decision.

This concern also has direct MAD precedent. A recent print--scan study reports substantial degradation of digitally trained detectors across acquisition devices~\cite{paulo2026printed}. Our storage and reconstruction controls ask a complementary question: can processing produce an attack response even when no second identity is present? This complements domain-shift evaluation with matched tests of score specificity.

\paragraph{Blending as supervision and as a control.}
Face X-ray learns a representation of blending boundaries for general face-forgery detection~\cite{li2020facexray}. Self-Blended Images synthesize forgery supervision by blending transformed versions of a single pristine image~\cite{shiohara2022sbi}. Both establish that useful manipulation supervision can be generated without examples from the target forgery method. Their objective is to recognize forgery traces; our same-person and seam-only negatives instead challenge the sufficiency of selected traces for a two-contributor label. A self-blend can therefore be a useful positive for forgery detection and a useful negative control for contributor dependence. This difference in target property, rather than a claim that blending artifacts are universally undesirable, motivates our construction.

Our contribution connects three tests: processing-only supervision as a benchmark counterexample, matched comparisons that separate rendering and blending from contributor changes, and a controlled study of adding matched negatives. The framework tests whether a proposed explanation survives a specified comparison. It neither introduces a new backbone nor certifies an identity mechanism; texture, noise averaging, geometry and the surround remain measurable alternative explanations.
